Soit $ABC$ un triangle dont les angles sont aigus et tel que $AB < AC$. Soit $\Omega$ son cercle
circonscrit, et $S$ le milieu de l'arc de cercle $\overset{\frown}{BC}$ contenant $A$. La perpendiculaire à
$(BC)$ passant par $A$ coupe le segment $[BS]$ en un point $D$, et recoupe le cercle $\Omega$ en un point
$E$ distinct de $A$. La parallèle à $(BC)$ passant par $D$ coupe la droite $(BE)$ en un point $L$. Enfin,
on note $\omega$ le cercle circonscrit au triangle $BDL$, et $P$ le point, autre que $B$, en lequel
$\omega$ recoupe le cercle $\Omega$.

Démontrer que la tangente à $\omega$ en $P$ et la droite $(BS)$ se coupent en un point situé sur la
bissectrice de $\widehat{BAC}$.
